\section{Partitions, shift averages, and controlled expansion}
\label{sec:geometry}

We continue with a nonempty common-base polytope
\(P=B(r_1)\cap B(r_2)\), and its center \(c\) from
Lemma~\ref{lem:center}.
The purpose of this section is to enlarge the two base polytopes in
carefully chosen directions while changing their common lattice-point
count by only a small relative error.  The enlargement will allow
small rank violations in the subsequent reduction.  Its proof uses
averages of counts with a small displacement between the two bases.

Three scales organize the comparison: \(S\) is a slack threshold for
the difference polytope, \(W\) is a larger slack threshold for the
individual base polytopes, and \(T\) is the amount of enlargement.
For the remaining estimates fix the integer scales
\begin{equation}\label{eq:scales}
 \begin{gathered}
 N=10^{10}+n+\ell+\lceil\eps^{-1}\rceil,
 \qquad K=N^5,\\
 S=N^{5000},\qquad W=N^{20000},\qquad T=W/N^{300}.
 \end{gathered}
\end{equation}
Every exponent used in this paper is an absolute constant.  Thus
quantities polynomial in $N$ are polynomial in the required size
parameters, although the degree is large.  We repeatedly use the
coarse bound $e^K$ for counts and box volumes.  More precisely,
\begin{equation}\label{eq:coarse-size-bound}
 \prod_{j=1}^m a_j\leq e^K
 \quad\text{if}\quad
 m\leq100n+100,\quad 0\leq a_j\leq N^{100000}(R+2).
\end{equation}
Indeed, when all factors are positive, the logarithm of the product
is at most
$(100n+100)(100000\log N+\log(R+2))<N^5$,
using $\log(R+2)\leq\ell+2$ and $N\geq10^{10}$.
Zero factors cause no difficulty.  When applying
\eqref{eq:coarse-size-bound} to a lattice count, the factors are the
numbers of allowed integer coordinate values; for volume they are
interval lengths.  These scales supply room for successive
approximations without any assumption on the widths of the input
polytope.  The precise bound proved in this section is an enlargement
factor of at most \(e^{20N^{-8}}\).

\subsection{Partitions determined by small slack}

For \(A\subseteq E\), define
\begin{equation}\label{eq:slack-functions}
 s_i(A)=r_i(A)-c(A),\qquad
 s_*(A)=s_1(A)+s_2(E\setminus A)
       =r_1(A)+r_2(E\setminus A)-R.
\end{equation}
These are nonnegative submodular functions, equal to zero at
\(\varnothing\) and \(E\).  By Lemma~\ref{lem:base-calculus},
\(B(s_*)=B(r_1)-B(r_2)\).

The atoms of a family of subsets of \(E\) are the nonempty classes of
elements with identical membership in every set of the family.  Let
\(\mathcal D_i\) be the partition into atoms of
\(\{A:s_i(A)<W\}\).  Let \(\mathcal G\) be the partition into atoms of
\(\{A:s_*(A)<S\}\); its parts will be called \emph{superblocks}.
We say that \(A\) cuts a part \(D\) if both \(A\cap D\) and
\(D\setminus A\) are nonempty.  Thus a set that cuts a
\(\mathcal D_i\)-part has \(s_i(A)\ge W\), and a set that cuts a
superblock has \(s_*(A)\ge S\).

We first record how bounds on a short generating family control the
totals on all its atoms.

\begin{lemma}[Prefixes and part totals]\label{lem:generator-prefixes}
Let \(\mathcal A=(A_1,\ldots,A_m)\) be a family of subsets of \(E\).
Its atoms admit an order such that every prefix is a union of at most
\(n\) intersections of members of \(\mathcal A\), allowing \(E\) as
the empty intersection.  If \(s\) is normalized, nonnegative, and submodular,
\(s(E)=0\), and \(s(A_j)\le b\) for every \(j\), where \(b\ge0\),
then in this order
every prefix \(F\) satisfies
\[
 s(F)\le n(m+1)b.
\]
Consequently, suppose that \(Q\subseteq B(r)\) and
\(r(A_j)-x(A_j)\le b\) for all \(x\in Q\) and all \(j\).
For every atom \(D\) of \(\mathcal A\) and every \(x,y\in Q\),
\begin{equation}\label{eq:part-range}
 |x(D)-y(D)|\le 2n(m+1)b.
\end{equation}
\end{lemma}

\begin{proof}
For an atom \(D\), intersect all generators containing \(D\), obtaining
\(H_D\).  This is the union of atoms whose membership sets contain
the membership set of \(D\).  Order the atoms so that a strict superset
of memberships comes first.  Every prefix is then the union of
\(H_D\) over its own atoms.  Nonnegative submodularity gives
\(s(U\cup V),s(U\cap V)\le s(U)+s(V)\), which proves the prefix
bound.  At any \(x\in Q\), apply that bound to the nonnegative slack
function \(r(\cdot)-x(\cdot)\).  For a fixed prefix its value lies
in \([0,n(m+1)b]\), so the prefix total varies by at most
\(n(m+1)b\).  Subtract consecutive prefix totals to obtain
\eqref{eq:part-range}.
\end{proof}

\begin{lemma}[Computable partitions and compatible orders]
\label{lem:partitions}
Each part of \(\mathcal D_i\) is contained in a superblock.
The three partitions can be found with polynomially many oracle calls
and bit operations.  Each has a generating family of at most \(n^2\)
sets satisfying its defining strict slack bound.

Fix a superblock order \(G_1,\ldots,G_g\) supplied by
Lemma~\ref{lem:generator-prefixes} for the \(s_*\)-generators.
Every prefix \(F_k=G_1\cup\cdots\cup G_k\), including the empty
and full prefixes, satisfies
\begin{equation}\label{eq:superblock-prefix-slack}
 s_*(F_k)\le N^4S.
\end{equation}
For system 1 use this order, and for system 2 use its reverse.  The
sets preceding a superblock \(G_k\) in these orders are therefore
\begin{equation}\label{eq:superblock-prefixes}
 A_{1G_k}=F_{k-1},\qquad A_{2G_k}=E\setminus F_k.
\end{equation}
\end{lemma}

\begin{proof}
If \(s_*(A)<S\), then \(s_1(A)<W\) and
\(s_2(E\setminus A)<W\), since \(S<W\).  Thus neither \(A\)
nor its complement cuts a fine part, proving the containment.

For each ordered pair of distinct elements \((a,b)\), minimize the
relevant slack function subject to \(a\in A\), \(b\notin A\).
Restricted submodular minimization, with rational modular terms for
\(s_i\), applies as in Section~\ref{sec:polyhedra}.  If the minimum
is below the threshold, retain a minimizer.  Two elements separable
by any set below the threshold are separated by a retained set in
one of the two orientations.  The retained family therefore has
exactly the required atoms, and contains at most \(n(n-1)\) sets.
Its values and all comparisons have polynomial encoding length.
The prefix bound is now Lemma~\ref{lem:generator-prefixes}, because
\(n(n^2+1)S\le N^4S\).
\end{proof}

Consider the bipartite multigraph with one vertex for every part of
\(\mathcal D_1\), one vertex for every part of \(\mathcal D_2\),
and one edge labeled \(j\) for each \(j\in E\), joining the two
parts containing \(j\).  Parallel edges are retained as distinct
coordinates.  Every connected component lies in a single superblock.
In each component \(C\), choose a representative coordinate
\(e_C\in E\).  In each superblock \(G\), choose a root component
\(C_0(G)\) and put \(j_G=e_{C_0(G)}\).  These choices require only
the computed partitions.  They will also be used in the port
parametrization of Section~\ref{sec:ports}.

\subsection{Two averages over small shifts}

For integral base polytopes \(B_1,B_2\) of total \(R\), put
\[
 Y_{B_1,B_2}(d)=
 \#\{(x_1,x_2)\in(B_1\cap\Z^E)\times(B_2\cap\Z^E):
                  x_1-x_2=d\}.
\]
Write \(Y(d)\) for the original pair, so \(Y(0)=Z\).
We shall also use the integral base polytopes
\begin{equation}\label{eq:expanded-bases}
 B_i^+=B(r_i)+\Delta_i,\qquad
 \Delta_i=
 \sum_{D\in\mathcal D_i}\ \sum_{\{a,b\}\subseteq D}
                 [-T,T](e_a-e_b),
\end{equation}
where each inner sum is over unordered two-element subsets.
The notation \(e_a\) in this formula denotes a unit vector, whereas
\(e_C\) denotes a chosen coordinate label.  The rank function of
\(B_i^+\) is
\begin{equation}\label{eq:expanded-rank}
 r_i^+(A)=r_i(A)+T\sum_{D\in\mathcal D_i}
                         |A\cap D|\,|D\setminus A|.
\end{equation}
Indeed, each segment has rank \(T\) on sets separating its two
endpoints and rank zero on other sets; use
Lemma~\ref{lem:base-calculus} for the sum.  In particular
\(r_i^+\) is integral and submodular, but need not be increasing.
Define \(Y^+(d)=Y_{B_1^+,B_2^+}(d)\) and \(Z^+=Y^+(0)\).

Every real vector with sum zero on each superblock has a unique
expression
\begin{equation}\label{eq:star-coordinates}
 d(l)=\sum_{G\in\mathcal G}\ \sum_{j\in G\setminus\{j_G\}}
                         l_j(e_j-e_{j_G});
 \qquad l_j=d_j.
\end{equation}
Call this the full set of star coordinates.  Its restricted version
uses only \(j=e_C\) for nonroot components \(C\subseteq G\),
with all other coefficients zero.  Set
\begin{equation}\label{eq:shift-kernels}
 h=S/N^{50},\qquad M=N^{2000},\qquad
 \omega(l)=\exp\left(-\frac{\sum_j|l_j|}{h}\right),\qquad
 \psi(l)=\exp\left(-\frac{\sum_C l_{e_C}^2}{M^2}\right).
\end{equation}
Here \(\omega\) uses the full star and \(\psi\) the restricted
star.  Let their unrestricted integer normalizers be
\begin{equation}\label{eq:shift-normalizers}
 A_\omega=\sum_{l\in\Z^{n-g}}\omega(l),\qquad
 J=\sum_{l\in\Z^{\sum_G(\#\{C\subseteq G\}-1)}}\psi(l).
\end{equation}
The sums factor into one-dimensional sums.  An empty star has a
single coefficient vector, weight one, and normalizer one.

\begin{lemma}[Shift averages]\label{lem:shift-average}
For \(F=Y\) and for \(F=Y^+\), every integral full or restricted
star shift with \(|l_j|\le S/N^{25}\) satisfies
\begin{equation}\label{eq:small-shift-ratio}
 e^{-N^{-8}}\le \frac{F(d(l))}{F(0)}\le e^{N^{-8}}.
\end{equation}
For \((\kappa,A_\kappa)=(\omega,A_\omega)\) or \((\psi,J)\),
on its corresponding star lattice,
\begin{equation}\label{eq:average-comparison}
 e^{-3N^{-8}}F(0)
 \le \frac{1}{A_\kappa}\sum_l\kappa(l)F(d(l))
 \le e^{3N^{-8}}F(0).
\end{equation}
The lower bound remains valid if the sum is restricted to
\(\max_j|l_j|\le S/N^{25}\).
\end{lemma}

\begin{proof}
First, \(Y(d)\ge1\) for every integral \(d\) with zero superblock
sums and \(\|d\|_1\le S/2\).  If a set \(A\) cuts no
superblock then \(d(A)=0\le s_*(A)\); otherwise
\(d(A)\le\|d\|_1\le S/2<s_*(A)\).  Thus \(d\in B(s_*)\),
so \(B(r_1)\cap(B(r_2)+d)\) is nonempty.  Its integrality,
from Lemma~\ref{lem:polyhedral-controls}, supplies an integer
point.  Inclusion gives the same positivity for \(Y^+\).

All coordinates of an expanded base lie in
\([-nT,R+nT]\).  The product-of-lengths bound
\eqref{eq:coarse-size-bound} therefore gives
\(Y(d),Y^+(d)\le e^K\) for every \(d\).  In the positive region
just proved, their logarithms lie between zero and \(K\).
By Proposition~\ref{prop:shift-logconcavity}, these logarithms
are concave on integer lines parallel to \(e_j-e_k\).

Introduce the star shifts successively.  Each partial shift has
\(\ell^1\)-norm at most \(2nS/N^{25}\).  On the line for the
next shift, an integer interval of radius \(\lfloor S/10\rfloor\)
about that partial shift remains in the positive region.  A concave
integer sequence taking values in \([0,K]\) on this interval has
successive slopes of absolute value at most \(100K/S\) within
distance \(S/N^{25}\) of its center: compare each slope with the
secants to the two ends of the interval.  Adding at most \(n\)
changes proves
\[
 |\log F(d(l))-\log F(0)|
 \le \frac{100K}{S}\sum_j|l_j|
 \le 100nK/N^{25}<N^{-8}.
\]

For completeness, put \(b=S/N^{25}\).  For an individual
\(\omega\)-coordinate the normalized tail beyond \(b\) is at
most \(4e^{-b/(2h)}\), by summing a geometric series.  For an
individual \(\psi\)-coordinate it is at most
\(4e^{-b^2/(2M^2)}\): compare the decreasing Gaussian tail with
its integral and use \(b\ge2M\).  The full normalized tail is
at most the sum of these one-coordinate tails.  Since
\(b/h=N^{25}\), \(b^2/M^2=N^{5950}\), and \(K=N^5\),
in either case the tail probability \(\tau\) satisfies
\[
 e^K\tau\le e^{-N^{20}}.
\]
On the complementary region use \eqref{eq:small-shift-ratio};
on the tail use \(F\le e^K\) and \(F(0)\ge1\).  The resulting
lower bound \((1-\tau)e^{-N^{-8}}F(0)\) and upper bound
\(e^{N^{-8}}F(0)+e^K\tau\) imply
\eqref{eq:average-comparison}, including its truncated version.
All assertions are exact when the star is empty.
\end{proof}

The full-star average is only a device for comparing \(Z^+\) and
\(Z\).  The algorithm will use the restricted average, whose shift
variables are determined by polynomial-length part totals.

\subsection{Subdivision into slices with slack}

The full-star numerator is a weighted sum over pairs of lattice
points with equal superblock totals.  Once these totals are fixed,
the two systems can be subdivided independently.  We will contract
them toward separate points in their respective slices; this
preserves the common superblock totals but may introduce differences
away from the representative coordinates.  The full star allows
these differences, which is why it is used for this comparison.
Fix such a vector of integer
totals \(w=(w_G:G\in\mathcal G)\), and write
\begin{equation}\label{eq:original-slices}
 Q_i(w)=\{x\in B(r_i):x(G)=w_G\ (G\in\mathcal G)\}.
\end{equation}
Whenever both expanded slices are nonempty, both original slices
are nonempty as well.  Every direction in \(\Delta_i\) preserves
the superblock totals, so the corresponding expanded slice is
exactly \(Q_i(w)+\Delta_i\).

The following refinement is an existence argument for these
slices.  Its subdivisions can differ between the two systems, and
can depend on \(w\).  Their number need not be polynomial; the
algorithm never enumerates them.

\begin{lemma}[Refinement witnesses]\label{lem:refinement-witnesses}
Fix integer \(w\) for which both \(Q_i(w)\) are nonempty.
For each system separately, its lattice points and those of its
expanded slice can be partitioned by fixing integer totals on
successively finer parts.  Each resulting original slice has the
form
\[
 Q=\{x\in B(r_i):x(D)=v_D\ (D\in\mathcal P)\},
\]
where \(\mathcal P\) refines \(\mathcal G\) and every part of
\(\mathcal P\) is a union of \(\mathcal D_i\)-parts.  Its
expanded partner is exactly \(Q+\Delta_i\).  Every nonempty
terminal slice has a point \(a^*\in Q\) satisfying
\begin{align}
 \|a^*-c\|_\infty&\le N^{60}S,
                       \label{eq:witness-distance}\\
 r_i(A)-a^*(A)&\ge S
       &&\text{if \(A\) cuts a part of \(\mathcal P\)},
                       \label{eq:witness-terminal-slack}\\
 r_i(A)-a^*(A)&\ge W/2
       &&\text{if \(A\) cuts a part of \(\mathcal D_i\)}.
                       \label{eq:witness-fine-slack}
\end{align}
\end{lemma}

\begin{proof}
We first control the starting totals.  If \(A\) is one of the
\(s_*\)-generators, it is a union of superblocks.  For any
\(x_1\in Q_1(w)\), \(x_2\in Q_2(w)\),
\[
 [r_1(A)-x_1(A)]+[r_2(E\setminus A)-x_2(E\setminus A)]
 =s_*(A)<S.
\]
Both terms are nonnegative.  Thus on the whole first slice every
generator has slack below \(S\), and on the whole second slice
its complement has slack below \(S\).  The same assertions hold
at \(c\).  Apply Lemma~\ref{lem:generator-prefixes} to
\(Q_i(w)\cup\{c\}\), using complemented generators for system 2.
It gives
\begin{equation}\label{eq:initial-total-distance}
 |w_G-c(G)|\le 2n(n^2+1)S\le N^6S.
\end{equation}
The distance bound in Lemma~\ref{lem:polyhedral-controls}, applied
to one rank system and these partition equalities, gives a point
\(a\in Q_i(w)\) with \(\|a-c\|_\infty\le N^{10}S\).
Choose such a point rationally: its allowed set is a nonempty
rational polytope, since \(c\) and the cube radius are rational.

We now describe one refinement step on one system.  Maintain a
nonempty original slice \(Q\), its partition \(\mathcal P\),
and a point \(a\in Q\).  A list of generators produces
\(\mathcal P\); initially it is the list just described.  The
list has at most \(n^2+n\) members.  Each initial generator has
slack at most \(S\) throughout the current slice, and every
subsequently added generator will have slack at most \(N^3S\)
throughout the slice on which it is added and all its descendants.

Intersect \(Q\) with the cube of radius \(L=N^{10}S\) about
\(a\), and choose a center \(b\) of this boxed slice with the
property of Lemma~\ref{lem:center}.  If every rank inequality
cutting a current part has slack at least \(S\) at \(b\), stop
on this branch with \(a^*=b\).  Otherwise select a set \(A\)
cutting a current part with \(r_i(A)-b(A)<S\).  The center
property implies
\[
 \max_{x\in Q,\ \|x-a\|_\infty\le L}
             (r_i(A)-x(A))
 \le 4(n+1)S<N^3S.
\]
The slack \(x\mapsto r_i(A)-x(A)\) is nonnegative on \(Q\)
and has integral linear part \(-\one_A\).  Since its displayed
maximum is less than \(N^3S<L\), the last assertion of
Lemma~\ref{lem:polyhedral-controls} identifies this maximum with
the unrestricted maximum on \(Q\).  Thus the new generator has
slack at most \(N^3S\) throughout the current slice.

We check simultaneously that the selected set does not cut a fine
part and that proximity to \(c\) survives refinement.  At depth
\(t\), maintain
\begin{equation}\label{eq:refinement-distance-induction}
 \|a-c\|_\infty\le N^{10}S+tN^{17}S.
\end{equation}
As each refinement increases the number of parts, \(t\le n\).
Every tested center \(b\), being within \(L\) of \(a\), is
therefore within \(N^{60}S\) of \(c\).  If \(A\) cut a fine
part, then its slack at \(c\) would satisfy
\(s_i(A)=r_i(A)-c(A)\ge W\), and so
\[
 r_i(A)-b(A)\ge W-nN^{60}S\ge W/2>S,
\]
a contradiction.  The refined partition obtained by adding
\(A\) still consists of unions of fine parts.

On this refined partition prescribe every integer vector of part
totals that occurs in the expanded current slice.  Every such
vector occurs in the original real slice: in a decomposition
\(x+v\in Q+\Delta_i\), the vector \(v\) preserves the new
totals, since it preserves each fine-part total.  Thus each
retained original child slice is nonempty, and its expanded
partner remains its Minkowski sum with \(\Delta_i\).  The new
generator list has slack at most \(N^3S\) throughout \(Q\).
By \eqref{eq:part-range}, the range of every new part total over
\(Q\) is at most
\[
 2n(n^2+n+1)N^3S\le N^{12}S.
\]
The current \(a\) consequently satisfies all the child's new
equalities with error at most \(N^{12}S\).  The distance bound
in Lemma~\ref{lem:polyhedral-controls} supplies a point in the
child within \(N^{17}S\) of \(a\).  As before, choose this point
rationally within the indicated rational cube.  This proves
\eqref{eq:refinement-distance-induction} at the next depth.

Each branch stops after at most \(n-1\) strict refinements.
There are finitely many branches because the slices are bounded.
The stopping rule gives \eqref{eq:witness-terminal-slack}; the
distance induction gives \eqref{eq:witness-distance}; and the
same comparison with \(c\) gives \eqref{eq:witness-fine-slack}
for every fine-cutting inequality.  Prescribing totals partitions
lattice points without multiplicity.  Integral prescribed totals
are compatible with the unit lattice in every terminal affine
space; in particular, nonempty original terminal slices contain
integer points by Lemma~\ref{lem:polyhedral-controls}.
\end{proof}

\subsection{From expanded cells to original cells}

We have produced points with two different slack guarantees.
The \(S\)-slack will absorb unit-cell rounding.  The larger
\(W/2\)-slack will also absorb the expansion in
\eqref{eq:expanded-bases}.  Both can be obtained by contracting
each slice by a tiny fraction toward its witness.

\begin{proposition}[Expansion costs little]\label{prop:expansion}
The expanded common-base count satisfies
\begin{equation}\label{eq:expansion-comparison}
 Z\le Z^+\le e^{20N^{-8}}Z.
\end{equation}
\end{proposition}

\begin{proof}
The lower bound follows from inclusion.  For the upper bound,
define the full-star weighted pair sums
\[
 \mathcal W=\sum_l\omega(l)Y(d(l)),\qquad
 \mathcal W^+=\sum_l\omega(l)Y^+(d(l)).
\]
Fix common integer superblock totals \(w\), and perform the two
independent subdivisions of Lemma~\ref{lem:refinement-witnesses}.
We compare one pair of terminal slices \(Q_1,Q_2\) and its
expanded pair \(Q_1+\Delta_1,Q_2+\Delta_2\).  Their partitions
\(\mathcal P_1,\mathcal P_2\) may differ.  Let their witnesses
be \(a_1^*,a_2^*\), and put \(a^*=(a_1^*,a_2^*)\).

For each system use the affine coordinates obtained by retaining
all but one physical coordinate in each terminal part.  The
omitted coordinate is the prescribed integer total minus the
retained coordinates.  The integral points in this affine space
therefore form precisely a unit grid.  Lebesgue measure below is
the product measure in these free coordinates, independently for
the two systems.  The slack in
\eqref{eq:witness-terminal-slack} ensures that each original
slice is full-dimensional in its specified affine space: every
inequality not fixed by its part totals is strict at the witness.
An affine space of dimension zero is given measure one.

Tile these affine spaces by half-open unit cells based at all
their integral points.  In physical coordinates any point of a
cell differs from its base point by at most \(2n\) in sup norm;
the same bound applies when rounding to the base of the containing
cell.  This uses the formula for the omitted coordinate and is
valid for negative coordinates as well.  No truncation at zero is
made: the expanded bases can have negative coordinates, and
feasibility will be checked using their full rank inequalities.

Let \(\mathcal C\) be the union of cells based at pairs of
original feasible lattice points in these terminal slices, and
let \(\mathcal C^+\) be the corresponding expanded union.
Extend the weight to their whole product affine space by
\[
 \omega(x_1,x_2)=e^{-\Phi(x_1,x_2)},\qquad
 \Phi(x_1,x_2)=\frac1h\sum_{G\in\mathcal G}
            \sum_{j\in G\setminus\{j_G\}}|x_{1j}-x_{2j}|.
\]
The common superblock totals ensure that this is the weight in
\eqref{eq:star-coordinates}.  It is continuous and \(\Phi\) is
convex.  Its logarithm varies by at most
\begin{equation}\label{eq:cell-weight-variation}
 \kappa_0=4n^2/h
\end{equation}
within any one of our product cells relative to the base point.
Thus integration over \(\mathcal C\) or \(\mathcal C^+\)
approximates the corresponding weighted lattice sum within a
factor \(e^{\kappa_0}\).

Put \(\theta=N^{-200}\) and apply the homothety
\[
 H(x)=(1-\theta)x+\theta a^*.
\]
We claim that \(H(\mathcal C^+)\subseteq\mathcal C\).
Fix \(x\in\mathcal C^+\), and round \(H(x)\) to the base
point \(y=(y_1,y_2)\) of its containing cell.  All prescribed
terminal totals are preserved.  In system \(i\), the rank
inequality for \(A\) is satisfied by the expanded lattice point
under \(x\) with additional allowance
\[
 \sigma_i(A)=T\sum_{D\in\mathcal D_i}
                         |A\cap D|\,|D\setminus A|
 \le n^2T.
\]
The move to \(x\) and the final rounding each change a subset
sum by at most \(2n^2\).  If \(A\) cuts a terminal part but
no fine part, then \(\sigma_i(A)=0\), and
\[
 y_i(A)-r_i(A)\le 4n^2-\theta S<0.
\]
If \(A\) cuts a fine part, it also cuts a terminal part, and
\[
 y_i(A)-r_i(A)
 \le 4n^2+n^2T-\theta W/2<0.
\]
These strict inequalities follow from the chosen scales:
\(\theta S=N^{4800}\) and
\(\theta W/T=N^{100}\).  Finally, if \(A\) cuts no terminal
part, its total is fixed on the affine space and equals a
feasible original total.  Every original rank inequality holds
at \(y_i\), proving the claim.  This also enforces nonnegativity
of \(y_i\), since it follows from the original base constraints.

Write \(d\le2n\) for the product free dimension.  The map \(H\)
is injective, with Jacobian \((1-\theta)^d\).  Convexity and
nonnegativity of \(\Phi\) give
\[
 \Phi(H(x))\le(1-\theta)\Phi(x)+\theta\Phi(a^*)
                  \le\Phi(x)+\theta\Phi(a^*).
\]
Moreover, \eqref{eq:witness-distance} gives
\[
 \theta\Phi(a^*)
 \le\theta\,2nN^{60}S/h<N^{-80}.
\]
Integration and the inclusion just proved imply
\[
 \int_{\mathcal C}\omega
 \ge (1-\theta)^d e^{-\theta\Phi(a^*)}
                         \int_{\mathcal C^+}\omega.
\]
Together with \eqref{eq:cell-weight-variation}, this bounds the
expanded weighted lattice sum in this pair of slices by the
original weighted sum times
\[
 \exp\{2\kappa_0+\theta\Phi(a^*)-d\log(1-\theta)\}
 \le e^{N^{-8}}.
\]
The same argument holds in dimension zero, with Jacobian and
cell measure one.  Summing over the terminal pairs and then
over \(w\) yields \(\mathcal W^+\le e^{N^{-8}}\mathcal W\).
Finally, Lemma~\ref{lem:shift-average} gives
\[
 Z^+\le e^{3N^{-8}}\frac{\mathcal W^+}{A_\omega}
 \le e^{4N^{-8}}\frac{\mathcal W}{A_\omega}
 \le e^{7N^{-8}}Z\le e^{20N^{-8}}Z.
\]
\end{proof}

Only the original partitions, their orders, and the representative
coordinates are computed.  The expanded polytopes, the full-star
sum, and all terminal subdivisions were used to prove the
comparison.  The next section uses the restricted Gaussian
average and part totals to construct a finite counting problem
whose capacities are polynomial in \(N\).
